
\section{Positionnement}

À RÉDIGER


% explication contrefactuelle locale, proximité et validité

% LTS formalisme de reflexion 

% Complexité qui explose -> evaluation relative à l'observation-> guider la recherche/generation si on sait evaluer la qualité

% integration de normes pour diriger notre evaluation evaluation experte, différentes normes pas de distinction

% modularité de la fonction permet d'obtenir de la diversite

% parser pddl

% Le formalisme utilisé dans ce travail pose l'hypothèse de Markov qui considère que l'état de départ d'une transition est le seul nécessaire pour déterminer celui d'arriver. De plus nous nous plaçons dans un cadre déterministe qui rend l'état d'arrivée unique.   
% Dans notre travail nous choisissons de rester dans le formalisme le plus général, nous restons au niveau des systèmes de transitions étiquetés
% Le choix du langages d'actions pour modeliser un système de transitions etiquetés dans le domaine du raisonnement sur l'action et le changement dépends de l'application, de l'expressivité souhaitée ou de la facilité de prise en main. Pour ne pas contraindre la recherche nous restons dans le formalisme le plus générale que sont les LTS. Même si ce cadre n'est pas utilisable de manière réaliste pour des applications complexes, rester dedans permet d'énoncer des resultats généraux pour tous les langages d'actions

% le travail n'est pas tourner vers la modélisation du monde on considère avoir une modèlisation qui contient tous les élements nécessaire et suffisant au rainsonnement. 

% nous on cherche à intégrer de la normativité dans les contrefactuelles (selon les théories morale c'est fait de manière un peu utilitariste) car : 